\chapter{METODE DAN RANCANGAN PENELITIAN}

\section{Deskripsi Umum Penelitian}

% TODO: Jelaskan gambaran besar penelitian yang akan dilakukan.

Lorem ipsum dolor sit amet, consectetur adipiscing elit. Sed do eiusmod tempor incididunt ut labore et dolore magna aliqua. Ut enim ad minim veniam, quis nostrud exercitation ullamco laboris nisi ut aliquip ex ea commodo consequat. Duis aute irure dolor in reprehenderit in voluptate velit esse cillum dolore eu fugiat nulla pariatur.

\section{Metodologi}

% TODO: Jelaskan metode penelitian yang digunakan.
% Contoh: Research and Development (R&D), Eksperimental, Studi Kasus, dll.

Tahapan penelitian dilakukan secara sistematis sebagai berikut:

\begin{enumerate}
    \item \textbf{Tahap 1:} Lorem ipsum dolor sit amet, consectetur adipiscing elit.
    \item \textbf{Tahap 2:} Sed do eiusmod tempor incididunt ut labore et dolore magna aliqua.
    \item \textbf{Tahap 3:} Ut enim ad minim veniam, quis nostrud exercitation ullamco laboris.
    \item \textbf{Tahap 4:} Duis aute irure dolor in reprehenderit in voluptate velit esse.
    \item \textbf{Tahap 5:} Excepteur sint occaecat cupidatat non proident.
\end{enumerate}

\section{Perancangan Sistem}

\subsection{Arsitektur Umum}

% TODO: Jelaskan komponen-komponen utama sistem yang dirancang.
% Sertakan diagram arsitektur jika ada.

Arsitektur sistem yang dirancang terdiri dari beberapa komponen utama:

\begin{itemize}
    \item \textbf{Komponen A:} Lorem ipsum dolor sit amet, consectetur adipiscing elit.
    \item \textbf{Komponen B:} Sed do eiusmod tempor incididunt ut labore et dolore magna aliqua.
    \item \textbf{Komponen C:} Ut enim ad minim veniam, quis nostrud exercitation ullamco laboris.
\end{itemize}

% TODO: Ganti dengan diagram arsitektur sistem Anda.
% \begin{figure}[H]
%   \centering
%   \includegraphics[width=0.9\textwidth]{figures/arsitektur-sistem}
%   \caption{Diagram Arsitektur Sistem}
%   \label{fig:arsitektur}
% \end{figure}

\subsection{Bagan Alir Metodologi}

% TODO: Tambahkan bagan alir (flowchart) metodologi penelitian.

% TODO: Ganti dengan bagan alir metodologi Anda.
% \begin{figure}[H]
%   \centering
%   \includegraphics[width=0.8\textwidth]{figures/bagan-alir}
%   \caption{Bagan Alir Metodologi Penelitian}
%   \label{fig:bagan-alir}
% \end{figure}

\section{Lingkungan Pengembangan}

% TODO: Jelaskan spesifikasi perangkat keras dan lunak yang digunakan.

Lingkungan pengembangan yang digunakan dalam penelitian ini adalah:

\begin{itemize}
    \item \textbf{Sistem Operasi:} Lorem ipsum dolor sit amet.
    \item \textbf{Perangkat Lunak:} Sed do eiusmod tempor incididunt.
    \item \textbf{Perangkat Keras:} Ut enim ad minim veniam.
\end{itemize}

\section{Metode Pengujian}

% TODO: Jelaskan metode pengujian yang akan digunakan.

Pengujian dilakukan dengan beberapa pendekatan:

\begin{enumerate}
    \item \textbf{Pengujian Fungsional:} Lorem ipsum dolor sit amet.
    \item \textbf{Pengujian Non-Fungsional:} Sed do eiusmod tempor incididunt.
    \item \textbf{Pengujian Lainnya:} Ut enim ad minim veniam.
\end{enumerate}
